\documentclass[a4paper, serif, 10pt]{moderncv}

%% moderncv
% style and color
\moderncvstyle{banking}
\moderncvcolor{black}

%% page layout/margins
\usepackage[top=2.5cm, bottom=2.5cm, left=1.5cm, right=1.5cm]{geometry}

\nopagenumbers{}

%% Babel config for Brazilian Portuguese language
% ref:
% - https://latex3.github.io/babel/guides/locale-portuguese.html
\usepackage[brazilian]{babel}

%% fontspec
\usepackage{fontspec}

\IfFontExistsTF{Linux Libertine}{
  \setmainfont[Ligatures={TeX, Common}, Numbers=OldStyle]{Linux Libertine}
}{}
\IfFontExistsTF{Linux Libertine O}{
  \setmainfont[Ligatures={TeX, Common}, Numbers=OldStyle]{Linux Libertine O}
}{}

%% CTeX
% CJK support, used to show CJK characters in the resume
%
% - fontset=none: disable builtin fontset but instead set the CJK font manually
% - heading=false: disable ctex heading
% - punct=kaiming: use kaiming punctuations styles for CJK
% - scheme=plain: use plain scheme, do not override \normalsize font size
% - space=auto: space settings for CJK characters
%
% ref:
% - http://ctan.mirrorcatalogs.com/language/chinese/ctex/ctex.pdf
\usepackage[UTF8, heading=false, punct=kaiming, scheme=plain, space=auto]{ctex}

\IfFontExistsTF{Noto Serif CJK SC}{
  \setCJKmainfont{Noto Serif CJK SC}
}{}
\IfFontExistsTF{Noto Sans CJK SC}{
  \setCJKsansfont{Noto Sans CJK SC}
}{}

%% line spacing
\usepackage{setspace}
\setstretch{1.125}

%% URL styling - use normal text instead of monospace
\urlstyle{same}

% Auto-underline all links
% Defer until \begin{document} so hyperref is already loaded
\AtBeginDocument{%
  \let\oldhref\href
  \renewcommand{\href}[2]{\oldhref{#1}{\underline{#2}}}%
}

%% Patch \httplink and \httpslink to handle full URLs
% The original moderncv macros always prepend http:// or https:// to URLs.
% This causes issues when the URL already has a protocol (e.g., https://example.com)
% resulting in malformed URLs like https://https://example.com.
% This patch checks if the URL already contains "://" and uses it directly if so.
% Note: We use \str_set:Nx (x-expansion) to expand macros like \@homepage first.
\ExplSyntaxOn
\str_new:N \g_yamlresume_url_str

\RenewDocumentCommand{\httpslink}{O{}m}{
  \str_set:Nx \g_yamlresume_url_str {#2}
  \str_if_in:NnTF \g_yamlresume_url_str {://}
    { \url{#2} }
    { \url{https://#2} }
}

\RenewDocumentCommand{\httplink}{O{}m}{
  \str_set:Nx \g_yamlresume_url_str {#2}
  \str_if_in:NnTF \g_yamlresume_url_str {://}
    { \url{#2} }
    { \url{http://#2} }
}
\ExplSyntaxOff

\name{Ana Beatriz Carvalho}{}
\title{Cientista de Dados Sênior com foco em Machine Learning e soluções analíticas}
\phone[mobile]{+55 11 98765-4321}
\email{ana.carvalho@datainsight.com.br}
\homepage{https://anabeatrizcarvalho.com.br}

\address{Av. Paulista, 1000, São Paulo, São Paulo, Brasil, 01310-100}{}{}

\extrainfo{{\small \faLinkedin}\ \href{https://www.linkedin.com/in/anabeatrizcarvalho}{@anabeatrizcarvalho} {} {} {} • {} {} {}
{\small \faGithub}\ \href{https://github.com/anabeatriz-data}{@anabeatriz-data} {} {} {} • {} {} {}
{\small \faMedium}\ \href{https://medium.com/@anabeatrizcarvalho}{@anabeatrizcarvalho}}

\begin{document}

\maketitle

\section{Informações Básicas}

\cvline{}{\begin{itemize}
\item Cientista de Dados com mais de 8 anos de experiência em modelagem preditiva, mineração de dados e inteligência de negócio.
\item Especialista em desenvolver soluções de machine learning para os setores financeiro e de varejo, gerando redução de custos e aumento de receita.
\item Sólida vivência em Python, SQL, computação em nuvem (AWS, GCP) e ferramentas de visualização de dados.
\item Experiência na liderança de times ágeis, na definição de indicadores e na comunicação de resultados para públicos não técnicos.
\end{itemize}}
\section{Formação}

\cventry{jan. de 2015–dez. de 2017}
        {Mestrado, Ciência de Dados, Nota: 9.1}
        {Universidade de São Paulo (USP)}
        {\url{https://www.usp.br/}}
        {}
        {\begin{itemize}
\item Dissertação sobre seleção de atributos e interpretabilidade de modelos preditivos.
\item Atuação como monitora nas disciplinas de estatística e aprendizado de máquina.
\end{itemize}
\textbf{Cursos}: Aprendizado de Máquina Avançado, Estatística Computacional, Mineração de Dados, Modelagem Preditiva, Visualização de Dados}

\cventry{fev. de 2010–dez. de 2014}
        {Bacharelado, Estatística, Nota: 8.7}
        {Universidade Estadual de Campinas (Unicamp)}
        {\url{https://www.unicamp.br/}}
        {}
        {\begin{itemize}
\item Iniciação científica em modelos de sobrevivência aplicados à saúde pública.
\item Participação em projetos de consultoria estatística para empresas locais.
\end{itemize}
\textbf{Cursos}: Inferência Estatística, Modelos Lineares, Amostragem, Análise Multivariada}
\section{Experiência Profissional}

\cventry{mar. de 2021–Atual}
        {Cientista de Dados Sênior}
        {DataWise Soluções em Dados}
        {\url{https://www.datawise.com.br}}
        {}
        {\begin{itemize}
\item Lidero o desenvolvimento de modelos de previsão de churn e propensão de compra para clientes de varejo e finanças.
\item Implementei um pipeline de MLOps que reduziu em 40\% o tempo de deploy dos modelos em produção.
\item Criei um framework de testes A/B e de validação offline que melhorou a confiabilidade dos experimentos.
\item Supervisiono cientistas de dados júnior e contribuo com a definição da arquitetura de dados.
\end{itemize}
\textbf{Palavras-chave}: Machine Learning, MLOps, Python, AWS, Airflow, Testes A/B}

\cventry{jan. de 2018–fev. de 2021}
        {Cientista de Dados}
        {VarejoTech}
        {\url{https://www.varejotech.com.br}}
        {}
        {\begin{itemize}
\item Construí modelos de precificação e segmentação de clientes que elevaram a margem operacional em 12\%.
\item Desenvolvi painéis analíticos em Power BI para acompanhamento de metas e detecção de anomalias.
\item Automatizei a coleta, limpeza e integração de dados de vendas e estoque com Python e SQL.
\end{itemize}
\textbf{Palavras-chave}: Segmentação, Preços, Power BI, SQL, Python, Previsão de Demanda}

\cventry{jan. de 2016–dez. de 2017}
        {Analista de Dados}
        {Grupo Financeiro Beta}
        {\url{https://www.grupobeta.com.br}}
        {}
        {\begin{itemize}
\item Atuei na construção de relatórios gerenciais de risco de crédito e inadimplência.
\item Realizei análises exploratórias para apoiar a área de crédito na definição de políticas de concessão.
\item Elaborei indicadores para comitês executivos e apresentações de resultados.
\end{itemize}
\textbf{Palavras-chave}: Risco de Crédito, Relatórios Gerenciais, SQL, Excel Avançado, Análise Exploratória}
\section{Idiomas}

\cvline{Português}{Nativo ou Bilíngue \hfill \textbf{Palavras-chave}: Português do Brasil}
\cvline{Inglês}{Proficiência Profissional Fluente \hfill \textbf{Palavras-chave}: TOEFL iBT 110}
\cvline{Espanhol}{Proficiência Profissional Limitada \hfill \textbf{Palavras-chave}: Leitura técnica}
\section{Competências}

\cvline{Machine Learning}{Especialista \hfill \textbf{Palavras-chave}: Python, Scikit-learn, TensorFlow, PyTorch, XGBoost, LightGBM, MLflow, Feature Store}
\cvline{Estatística}{Especialista \hfill \textbf{Palavras-chave}: Inferência Bayesiana, Séries Temporais, Testes A/B, Modelos Lineares, Modelos de Sobrevivência}
\cvline{Engenharia de Dados}{Avançado \hfill \textbf{Palavras-chave}: SQL, Apache Spark, Airflow, dbt, Docker, AWS, GCP}
\cvline{Visualização de Dados}{Avançado \hfill \textbf{Palavras-chave}: Matplotlib, Seaborn, Plotly, Power BI, Looker}
\section{Prêmios}

\cventry{nov. de 2023}
        {Associação Brasileira de Análise de Dados}
        {Prêmio ABRADe de Inovação em Dados}
        {}
        {}
        {Reconhecimento pelo projeto de MLOps aplicado à prevenção de churn em uma base com mais de 5 milhões de clientes.}

\cventry{dez. de 2014}
        {Universidade Estadual de Campinas}
        {Destaque Acadêmico}
        {}
        {}
        {Reconhecimento pela excelência no desempenho do curso de Estatística entre os concluintes daquele ano.}
\section{Certificados}

\cventry{jun. de 2022}
        {Amazon Web Services}
        {AWS Certified Machine Learning - Specialty}
        {\url{https://aws.amazon.com/certification/}}
        {}
        {}

\cventry{mar. de 2021}
        {Google}
        {TensorFlow Developer Certificate}
        {\url{https://www.tensorflow.org/certificate}}
        {}
        {}
\section{Publicações}

\cventry{set. de 2020}
        {Previsão de inadimplência em carteiras de crédito usando aprendizado de máquina}
        {Revista Brasileira de Inteligência Computacional}
        {\url{https://www.rbic.org.br}}
        {}
        {\begin{itemize}
\item Apresenta uma comparação entre regressão logística, florestas aleatórias e redes neurais para predição de inadimplência.
\item Discute estratégias de balanceamento de classes e interpretabilidade de modelos em cenários reais de crédito.
\end{itemize}}
\section{Referências}

\cventry{\emaillink[ricardo.nogueira@datawise.com.br]{ricardo.nogueira@datawise.com.br}}
        {Gerente de Dados na DataWise Soluções em Dados}
        {Ricardo Nogueira}
        {+55 11 99876-5432}
        {}
        {Ana Beatriz combina sólida formação estatística com capacidade de liderança e comunicação. Foi fundamental para a evolução dos nossos modelos de machine learning.}
\section{Projetos}

\cventry{mar. de 2022–Atual}
        {Pipeline end-to-end que combina automação de treinamento, validação e deploy de modelos de churn.}
        {Plataforma de MLOps para Previsão de Churn}
        {\url{https://github.com/anabeatriz-data/churn-mlops}}
        {}
        {\begin{itemize}
\item Utiliza Docker, Airflow e MLflow para orquestrar as etapas do ciclo de vida do modelo.
\item Reduziu o tempo de atualização dos modelos de semanas para horas no ambiente de produção.
\end{itemize}
\textbf{Palavras-chave}: MLOps, Docker, Airflow, MLflow, Churn}

\cventry{jan. de 2020–dez. de 2020}
        {Projeto de estudo com modelos ARIMA, Prophet e redes LSTM aplicados a uma base pública de vendas.}
        {Modelo de Séries Temporais para Previsão de Vendas}
        {\url{https://github.com/anabeatriz-data/series-temporais-vendas}}
        {}
        {\begin{itemize}
\item Análise comparativa de acurácia entre abordagens clássicas e neurais.
\item Desenvolvimento de um dashboard interativo em Streamlit para exploração dos resultados.
\end{itemize}
\textbf{Palavras-chave}: Séries Temporais, ARIMA, Prophet, LSTM, Streamlit}

\cventry{ago. de 2019–out. de 2019}
        {Dashboard criado com Power BI para monitoramento de metas comerciais e qualidade dos dados.}
        {Painel de Indicadores de Negócio}
        {\url{https://github.com/anabeatriz-data/dashboard-indicadores}}
        {}
        {\begin{itemize}
\item Consolidou mais de 15 fontes de dados em um único painel gerencial.
\item Facilitou a tomada de decisão semanal dos times comerciais.
\end{itemize}
\textbf{Palavras-chave}: Power BI, Indicadores, DAX, SQL}
\section{Interesses}

\cvline{Leitura}{Ficção científica, Não ficção}
\cvline{Esportes}{Corrida, Yoga}
\section{Voluntariado}

\cventry{jan. de 2020–Atual}
        {Mentora voluntária}
        {Mulheres em Dados}
        {\url{https://www.mulheresemdados.org.br}}
        {}
        {\begin{itemize}
\item Mentoria de carreira para mulheres que estão ingressando na área de dados.
\item Criação de conteúdos sobre machine learning e estatística para a comunidade.
\end{itemize}}


\cventry{mar. de 2019–dez. de 2022}
        {Voluntária em projetos sociais}
        {Ciência de Dados para o Bem}
        {\url{https://www.datascienceforgood.org.br}}
        {}
        {\begin{itemize}
\item Atuação em projetos de análise de dados para organizações sem fins lucrativos nas áreas de educação e saúde.
\end{itemize}}

\end{document}